\begin{activity}{Linear or Non-linear?}

\section*{Purpose}

This activity builds the skill of classifying geophysical forward relations as linear, non-linear, or linearizable, using a range of equations drawn from gravity, magnetics, geothermics, and physical property estimation. Correct classification determines whether a problem can be solved in a single step or requires iterative inversion.

\section*{Learning Outcomes}

By the end of this activity, you should be able to:
\begin{itemize}
  \item Identify the data, model parameters, and operator in a geophysical forward relation.
  \item Apply the homogeneity and additivity tests to determine whether a relation is linear in its model parameters.
  \item Propose a linearization strategy (log-transform, change of variables, or Taylor expansion) for non-linear relations where one exists.
  \item Explain the practical consequences of linearity vs.\ non-linearity for inversion.
\end{itemize}

\section{Theory}

Inversion problems can be cast as \(\mathbb{A}\mathbf{m} = \mathbf{d}\) for linear problems and \(A(m) = d\) for non-linear problems.  Where \(d\) is the data and \(A\) is the operator or design (more on that below).  The difference may appear subtle, though in reality the difference is that the linear equation can be written as a linear algebra problem as the model parameters are independent and discrete.  In the non-linear formulation, the one or more model parameters are dependent upon another.  As a result they cannot be fully separated unless they can be linearized.

Let's explore what \(A\), \(m\), and \(d\) are a bit more intuitively.
The data, \(d\), are the results of an experiment.  Identifying the data are generally straightforward.  The operator and model parameters can takes a bit more getting used to.  The model parameters, \(m\), are the physical or virtual properties of the system and initially unknown.  Once you determine these properties, the forward model becomes a prediction.  The operator or design matrix, \(A\), can be thought of \emph{as the experiment itself}, and is dependent upon the geometry, conditions of the experiment, and the physics.  The model tells us what exists, the operator tells us how we observe it.

When we are doing a geophysical experiment with the Earth as the medium,
\begin{itemize}
  \item \(\mathbf{m} \to \) the model is Earth, what it is made of and its physical properties. Changing \(\mathbf{m}\) changes the Earth.
  \item \(\mathbf{A} \to \) how Earth talks to the instrument, how the Earth is sampled by the experiment, i.e., the rules of the experiment. Changing \(\mathbf{A}\) changes the experiment, i.e., same Earth + different experiment \(\to\) different data.
  \item \(\mathbf{d} \to \) the data is the outcome of the experiment, what we measure
\end{itemize}

In most geophysical problems, \(\mathbf{A}\) is
\begin{itemize}
  \item gravity -- distances from mass elements
  \item magnetics -- sensor postions and dip angle
  \item heat flow -- layer thicknesses and boundary conditions
  \item electromagnetics -- distances from conductive/resistive bodies
  \item seismology -- ray paths or travel times
\end{itemize}
So the entries in \(\mathbf{A}\) are determined by:
\begin{itemize}
  \item positions
  \item distances
  \item times
  \item physical laws
  \item boundary conditions
\end{itemize}

Put simply, the matrix \(\mathbb{A}\) encodes the physics and geometry of how the experiment samples the model, \(\mathbf{m}\), resulting in the observations, \(\mathbf{d}\).

\section*{Linearizing}

Many non-linear equations can be transformed into a linear form. Doing so is a great advantage because it allows for the problem to be solved using linear least squares methods and the solution is obtained in one step. Before resorting to a non-linear inversion, ask whether one of the following approaches applies.

\subsection*{Logarithms for power laws}
If $y=ax^b$ take the logarithm:
\[
  \log y = \log a + b \log
\]
which is linear in $\ln x$.  Examples: Gardner's relation, many petrophysical correlations,
empirical scaling laws.

\subsection*{Logarithms for exponentials}
If $y=ae^{\pm bx}$ then
\[
  \log y = \log a \pm bx
\]
which is linear.  Examples: radioactive decay, compaction models, and attenuation relations.

\subsection*{Change of variables}
Sometimes a product of unknowns behaves as a single parameter (i.e., are inseperable).  For example from the gravity of a buried sphere, $g \propto a^3\Delta\rho$.  Define:
\[
  M=a^3\Delta\rho
\]
then $g \propto M$, which is linear in the new parameter.  This is particularly useful for demonstrating why some parameters are not independently resolvable.

\subsection*{Reciprocal transforms}
Some rational functions become linear after taking an inverse:
\[
  y=\frac{1}{a+bx}
\]
gives
\[
  \frac{1}{y}=a+bx
\]
which is linear.

{\emph Note: When transforming, as for many of the examples above, the data, model parameters and observation conditions all have to be transformed in order to achieve a linear form.  Once the least squares solution is solved, the parameters can be transformed into normal space to express the model in its original form.}

\section*{Task}

For each of the following forward relations drawn from potential fields, geothermics, and physical property relations, decide whether it is (i) linear, (ii) non-linear, or (iii) non-linear but can be linearized. If non-linear, describe how you would linearize or re-parameterize it. Justify your choice with respect to the model parameters that would be inverted.

In your groups:
\begin{enumerate}[label=\arabic*.,itemsep=2pt]
  \item Identify the data and the experimental quantities.
  \item For each equation
  \begin{itemize}
    \item identify the quantity or quantities that would be treated as the \emph{model parameter(s)} in an inversion.
    \item Classify the mapping data $\to$ model as linear or non-linear in those parameters.
    \item If non-linear, propose a linearization (e.g., log-transform, Taylor expansion/Jacobian around $p_0$, or a change of variables) and note any loss of identifiability.
  \end{itemize}
\end{enumerate}

\subsection*{Equations}
\begin{enumerate}[label=Q\arabic*.,itemsep=6pt]
  \item \textbf{Gravity (Bouguer slab):}
  We'll see this equation again as it is used to model the gravity due to an infinite slab.
  \[
    \Delta g = 2\pi G(\rho_s - \rho_c)H
  \]
  \answerbox[2cm]

  \item \textbf{Gravity (buried sphere):} 
  This gravity solution is one of the simplest models for a gravity anomaly.  We will explore it further.
  \[
    g(r) = \tfrac{4}{3}\, \pi a^3 G\, (\rho_s - \rho_c) \, \dfrac{z}{(r^2+z^2)^{3/2}}
  \]
  \answerbox[2cm]

  \item \textbf{Thermal conductivity $k(T,X)$:} 
  We'll use this formulation in the problem set.  This equation is used  to estiamte thermal conductivity model for a solid solution series for a mineral.
  \[
    k(T,X) = (a + bX + cX^2)\,(298/T)^n
  \]
  \answerbox[2cm]

  \item \textbf{Thermal conductivity $k(T,P)$:} 
  This form of thermal conductivity is often used to estimate thermal conductivity for a rock.
  \[
    k(T,P) = \dfrac{k_0}{A + BT}\,(1 + cP)
  \]
  \answerbox[2cm]

  \item \textbf{Isostasy (Airy):}
  Airy isostasy is an end-member isostatic model, used to explain topographic variations due to differences in crustal thickness.
  \[
    \Delta t = \dfrac{\rho_c}{\rho_m - \rho_c}\,h
  \]
  \answerbox[2cm]

  \item \textbf{Isostatic gravity correction (column):} 
  Gravity can be corrected for a mass deficit created by variations in crustal thickness via an isostatic correction.
  \[
    \Delta g_{\mathrm{iso}} \approx 2\pi G(\rho_m-\rho_c)\,\Delta t - 2\pi G\,\rho_c\,h
  \]
  \answerbox[2cm]

  \item \textbf{Magnetics (2-D infinite dike):}
  The ocean floor records chaotic reversals in Earth's magnetic field.  The remnant magnetic field is ``frozen'' in dikes produed at the mid-ocean ridge.
  \[
    \Delta T(x) = C\Big[\arctan\Big(\frac{x-x_1}{z}\Big) - \arctan\Big(\frac{x-x_2}{z}\Big)\Big]
  \]
  \answerbox[2cm]

  \item \textbf{Layered 1-D heat with sources:}
  Geotherms can largely be explained by 1-D heat flow through a layered Earth.  We'll see this again when we start geothermics.
  \[
    T_i(z) = T_{i-1} + \frac{q_{i-1}}{k_i}(z_i - z_{i-1}) - \frac{A_i}{2k_i}(z_i - z_{i-1})^2
  \]
  \answerbox[2cm]

  \item \textbf{Density--porosity mix:}
  Sedimentary rocks include pore space (porosity) that may is filled with a fluid (air, water, brine, gas, oil, etc.).
  \[
    \rho_\mathit{bulk} = (1-\phi)\,\rho_\mathit{matrix} + \phi\,\rho_\mathit{fluid}
  \]
  \answerbox[2cm]

  \clearpage
  \item \textbf{Gardner's relation (density--velocity):}
  Many solutions in geophysics involve trying to estimate one property from another, in this case density from seismic velocity.
  \[
    \rho = a\,v^b
  \]
  \answerbox[2cm]

  \item \textbf{Gravity (discrete 2-D prisms):}
  While we can produce analytic equations for a variety of idealized geologic bodies with simple geometries (e.g., buried spheres as above), to produce more geologically realistic bodies we can produce structures from arbitrarily-shaped polygons.
  \[
    g_z =
    2G\,\Delta\rho
    \sum_{i=1}^{N}
    \left[
    (x_{i+1}-x_i)\log\!\left(\frac{r_{i+1}}{r_i}\right)
    +
    (z_{i+1}-z_i)
    \arctan\!\left(
    \frac{x_i z_{i+1} - x_{i+1} z_i}
    {x_i x_{i+1} + z_i z_{i+1}}
    \right)
    \right]
  \] 
  where
  \[
    r_i = \sqrt{x_i^2 + z_i^2}, 
    \qquad
    (x_{N+1}, z_{N+1}) = (x_1, z_1).
  \]
  \answerbox[2cm]
\end{enumerate}

\section*{Key Ideas}

\begin{minipage}[t]{0.58\textwidth}
\textbf{The operator $\mathbf{A}$ encodes the physics and geometry of the experiment.} Changing the Earth (the model $\mathbf{m}$) changes the predictions; changing the experiment (the operator $\mathbf{A}$) also changes the predictions even if the Earth stays the same.

\vspace{0.5em}
\textbf{Many geophysical relations are non-linear but can be linearized.} Power laws become linear under a log-transform; functions with products of unknowns may become linear under a change of variables. Recognising these opportunities converts a difficult iterative problem into a tractable one.

\vspace{0.5em}
\textbf{Non-linearity has consequences for inversion}: the Jacobian must be recomputed at each iteration, convergence is not guaranteed, and the solution found may be a local rather than a global minimum. The choice of starting model matters far more for non-linear than for linear problems.
\end{minipage}
\hfill
\begin{minipage}[t]{0.38\textwidth}
    \begin{tikzpicture}[
            >=latex,
            node distance=1.0cm,
            process/.style={
                rectangle,
                rounded corners,
                draw,
                align=center,
                minimum width=2.5cm,
                minimum height=0.8cm,
                font=\small
            },
            decision/.style={
                diamond,
                draw,
                aspect=2,
                minimum width=2.5cm,
                align=center,
                inner sep=1pt,
                font=\small
            },
            baseline=(current bounding box.north)
        ]

        \node[process] (model)
            {Forward\\Model};

        \node[decision, below=of model] (linear)
            {Linear?};

        \node[process, below=of linear] (llsq)
            {Linear Least\\Squares};

        \node[decision, below right=3.5cm and 2.5cm of linear] (linz)
            {Can be\\Linearized?};

        \node[process, below=1.3cm of llsq] (transform)
            {Transform\\to Linear};

        \node[process, below=of linz] (nonlinear)
            {Non-linear\\Inversion};

        % Main flow
        \draw[->] (model) -- (linear);

        \draw[->] (linear) -- node[left] {Yes} (llsq);

        \draw[->]
            (linear.east)
            -| node[pos=0.25,above] {No}
            (linz.north);

        \draw[->] (linz) -- node[right] {No} (nonlinear);

        \draw[->] (linz.west) -- node[midway, below] {Yes} (transform);

        \draw[->] (transform) -- (llsq);
    \end{tikzpicture}
\end{minipage}

\end{activity}